% T32 — «Финал года»: официальная итоговая работа, синяя шапка с гербом,
% части А (тест) и Б (задачи), нумерация А1, А2, Б1. Math, класс 9.
\documentclass[a4paper,11pt]{article}

\usepackage{../../../../Lessons/_templates/preamble_worksheet}
\usepackage{../_shared/worksheet_check}

% --- Палитра: официальный тёмно-синий + золото ---
\definecolor{wcprimary}{HTML}{0E2A56}
\definecolor{wcprimaryLight}{HTML}{B5C3DD}
\definecolor{wcprimaryBG}{HTML}{F2F4FA}
\definecolor{wcgold}{HTML}{C9A227}

% Шрифт оставляем serif (Times) для официальности

% --- Заголовок с гербом-щитом ---
\renewcommand{\WorksheetTitle}[1]{%
  \par\noindent\begin{tikzpicture}
    \fill[wcprimary] (0,0) rectangle (\linewidth, -22mm);
    % Тонкая золотая полоса
    \fill[wcgold] (0,-22mm) rectangle (\linewidth, -23mm);
    % Щит-герб (упрощённый)
    \begin{scope}[shift={(15mm, -11mm)}, scale=0.6]
      \fill[wcgold] (0,0) -- (4mm, 0) -- (4mm, -5mm) .. controls (4mm,-9mm) and (0,-10mm) .. (0,-12mm) -- cycle;
      \fill[wcgold] (0,0) -- (-4mm,0) -- (-4mm,-5mm) .. controls (-4mm,-9mm) and (0,-10mm) .. (0,-12mm) -- cycle;
      \draw[white, line width=0.5pt] (0,0) -- (0,-12mm);
      % Звезда внутри
      \foreach \a in {90, 162, 234, 306, 18}
        \fill[wcprimary] (\a:1.5mm) -- (\a+36:0.6mm) -- (\a-36:0.6mm) -- cycle;
    \end{scope}
    \node[anchor=west, text=white, font=\bfseries\Large]
      at (25mm, -7mm) {ИТОГОВАЯ РАБОТА};
    \node[anchor=west, text=wcgold, font=\bfseries]
      at (25mm, -13mm) {за учебный год};
    \node[anchor=west, text=wcprimaryLight, font=\itshape]
      at (25mm, -18mm) {#1};
  \end{tikzpicture}\par\vspace{5mm}}

% Разделитель частей
\newcommand{\PartBanner}[2]{% #1=буква, #2=название
  \par\vspace{4mm}\noindent\begin{tikzpicture}
    \fill[wcprimary] (0,0) rectangle (\linewidth, -8mm);
    \fill[wcgold] (0,-8mm) rectangle (\linewidth,-8.6mm);
    \node[anchor=west, text=wcgold, font=\bfseries\Large]
      at (3mm, -4mm) {ЧАСТЬ\,#1};
    \node[anchor=west, text=white, font=\bfseries]
      at (28mm, -4mm) {#2};
  \end{tikzpicture}\par\vspace{3mm}}

% TaskBox с префиксом-буквой части (А1, А2, Б1, ...)
\newcommand{\ExamTask}[3]{% #1 префикс (А или Б) #2 номер #3 содержимое
  \par\noindent\begin{tikzpicture}
    \node[draw=wcprimary, line width=0.7pt, rounded corners=2pt,
          fill=white, inner sep=3mm,
          text width=\dimexpr\linewidth-6mm\relax,
          anchor=north west] (B) at (0,0) {#3};
    \node[anchor=south west, fill=wcgold, text=wcprimary,
          rounded corners=1pt, inner xsep=5pt, inner ysep=2pt,
          xshift=1.5pt] at (B.north west)
      {\scriptsize\bfseries Задание #1#2};
  \end{tikzpicture}\par\vspace{2.5mm}}

\SetAnswerFieldSize{30mm}{8mm}

\begin{document}

\WorksheetTitle{Математика. 9 класс}
{\small\itshape\color{wcprimary!85}T32 \quad·\quad ID: T32-final-2026-05-27 \quad·\quad на выполнение~--- 90 минут}\par\vspace{2mm}
\FirstPageHeader

\PartBanner{А}{краткий ответ}

\ExamTask{А}{1}{%
  Найдите значение выражения $\dfrac{2{,}8 \cdot 1{,}5}{0{,}7}$.\par
  \vspace{1mm}\hfill\textbf{\color{wcprimary}Ответ:}~\AnswerField{1}%
}

\ExamTask{А}{2}{%
  Решите уравнение $x^2 - 9x + 14 = 0$. В ответе укажите больший корень.\par
  \vspace{1mm}\hfill\textbf{\color{wcprimary}Ответ:}~\AnswerField{2}%
}

\ExamTask{А}{3}{%
  В арифметической прогрессии $a_1 = 7$, $d = 4$. Найдите $a_{12}$.\par
  \vspace{1mm}\hfill\textbf{\color{wcprimary}Ответ:}~\AnswerField{3}%
}

\ExamTask{А}{4}{%
  Найдите площадь прямоугольного треугольника с катетами $8$~см и $15$~см.\par
  \vspace{1mm}\hfill\textbf{\color{wcprimary}Ответ:}~\AnswerField{4}%
}

\PartBanner{Б}{задачи с подробным решением}

\ExamTask{Б}{1}{%
  Велосипедист проехал из пункта~А в пункт~Б за $4$ часа, а обратно~--- за $5$ часов,
  так как скорость уменьшилась на $3$~км/ч. Найдите расстояние между~А и~Б.\par
  \vspace{1mm}\hfill\textbf{\color{wcprimary}Ответ (км):}~\AnswerField{5}%
}

\ExamTask{Б}{2}{%
  В прямоугольном треугольнике гипотенуза равна $25$, а один из катетов равен $7$.
  Найдите площадь треугольника.\par
  \vspace{1mm}\hfill\textbf{\color{wcprimary}Ответ:}~\AnswerField{6}%
}

\WorksheetQR{T32-final/L1/v1}

\end{document}
